\begin{frame}
	\ \\
	\noindent В 2019 году олимпиада «Математика НОН-СТОП» проводится
	с использованием гранта Президента Российской Федерации \\
	на развитие гражданского общества, предоставленного \\
	{\bfseries Фондом президентских грантов.}
	
	\vspace{0.45cm}
	\begin{center}
		\includegraphics[scale=0.6]{fpg}
	\end{center}
\end{frame}

\begin{frame}\frametitle{История олимпиады}
\begin{itemize}
	\item 2010 — первая олимпиада;
        %\mitem 2011 — 418 участников
%	\mitem 2012 — три площадки;
%	\mitem 2015 — первый профильный вариант;
	\item 2016 — 400 участников пишут базовый вариант, 92 --- профильный;\\
        $\phantom{2016}$ — поддержка Фонда «Время Науки»;
	%\mitem 2017 — задачи для 4 класса;\\
	%$\phantom{2017}$ — автоматическое распределение участников;
	\item 2018 — 847 участников пишут базовый вариант, 128 --- профильный;\\
        $\phantom{2018}$ — включение в Перечень региональных олимпиад и конкурсов\\
	$\phantom{2018 — }$\quad интеллектуальной направленности;\\
	$\phantom{2018}$ — поддержка Фонда Президентских грантов,\\
	$\phantom{2018 — }$\quad Комитета по образованию СПб;\\
	%$\phantom{2018}$ — этот семинар;\\
	\item 2019 — выход сборника задач;\\
        $\phantom{2019}$ — площадки в Бердске (Новосибирская обл.) и Гомеле (Беларусь);\\
        $\phantom{2019}$ — число участников приближается к 2000.
\end{itemize}\end{frame}

\begin{frame}\frametitle{Статистика олимпиады}
\begin{itemize}
	\item 12 площадок (на 2021 год — 20 соглашений);\\
	\item количество участников --- около 2000;\\
	\item пять городов: Санкт-Петербург, Бердск (Новосибирская обл.), \\
		Нов. Уренгой (ЯНАО), Гатчина (ЛО), \\
		Гомель (Беларусь);\\
	\item две страны;\\
	\item проблемы с часовыми поясами.
\end{itemize}\end{frame}

\begin{frame}\frametitle{Не верьте им!}
	Основная цель олимпиады — избегать появления «сообщества профессионалов», которые умеют её решать. \ps
	То, что вам сейчас расскажут, — не руководство по подготовке и не классификация задач будущей олимпиады; \ps
	А набор наблюдений, касающихся заданий прошлых лет.
\end{frame}
